\section{Love Gives Being}

The ascent began with the creature's confession: I have not given being to myself. The Thomistic account of divine love, now read in the light of revelation, permits a stronger statement without turning creation into necessity: the creature exists because God freely wills its good.

Aquinas marks the difference between created and divine love. Human love ordinarily encounters or imagines goodness and is moved by it. The goodness of the beloved is not caused simply by being loved. God's will, however, is the cause of created being and every created perfection. To love is to will good to another; therefore God's love does not find an independently existing goodness that persuades him. It gives and sustains the goodness by which the creature is lovable. In Aquinas's compressed phrase, divine love ``infuses and creates goodness in things'' (ST I, q.~20, a.~2).

The Book of Wisdom supplies the scriptural resonance: God loves all things that exist, detests nothing he has made, and preserves what he has called into being (Wis 11:24--26). This is not because nonbeing exerted a claim or because the world forced a choice. The gratuitousness remains absolute. Love is not an explanation that makes creation necessary; it names the free willing of the creature's good by the One who needs nothing from it.

Every created perfection can now be read as gift without becoming a coded message whose details we confidently decipher. The existence of a sparrow, a distant galaxy, a difficult neighbor, or the person reading is not self-grounded. Each is known and willed in the one divine act. Not everything that happens within creation is morally good, and the mystery of evil requires distinctions this article cannot supply in full. But whatever positive actuality a creature possesses is not outside the creative love that causes being and goodness; evil as evil is privation, not a rival substance produced by another first principle.

\subsection{The ascent is answered by a descent}

Metaphysics climbs from effects toward their cause. The Gospel according to John announces the reverse movement: the Word through whom the world is made enters the world. The First Letter of John does not leave divine love at the level of invisible causality. Love is manifested because the Father sends the only Son into the world that we might live through him; its measure is not our initiative but God's prior self-gift (1 John 4:9--10, 14).

The Cross therefore interprets the sentence ``God is love.'' It does not cause the Trinity to become loving. It manifests eternal divine charity in history; the missions of Son and Spirit express their eternal processions (ST I, q.~43, aa.~1--2). The one who suffers is truly the divine Son in the human nature he assumed; the divine nature does not become mutable or passible (CCC 468--469). Christian love is neither a cosmic emotion nor the romanticization of pain. It is the self-giving holiness by which God seeks the creature's true good at the cost borne in the flesh of the Word.

St.~Athanasius sees creation and redemption united in the one Word. The Word who called humanity from nonbeing is the Word who assumes a body to renew the corrupted image and conquer death. Only the maker can recreate; only Life can restore the mortal to incorruption (\emph{On the Incarnation} 4--10, 20, 54). Ontology opens into salvation without becoming identical with it. Natural existence is gift; life in Christ is a further, supernatural participation by grace.

\subsection{Love returning to Love}

Created love cannot equal its source. It can receive and answer. St.~Bernard of Clairvaux describes the Bride's love as returning to its principle: the little stream cannot match the fountain, yet what flows back was first received from the fountain (\emph{Sermons on the Song of Songs} 83.3--6). St.~John of the Cross presses the mystery into Trinitarian participation: in union by grace the soul loves God with the love God communicates, without ceasing to be a creature (\emph{Spiritual Canticle}, stanza 39.4--6, second redaction).

This is neither absorption nor equality of essence. The distinction established by creation remains forever. Grace does not make the creature uncreated; it gives a created participation in divine life (2 Pet 1:4). The more completely the creature receives, the more completely it becomes itself. Union is not erasure but communion.

John tests contemplation by love of the brother. Whoever claims to love the unseen God while refusing the neighbor contradicts the revealed form of participation (1 John 4:20--21). Augustine makes the same movement in \emph{De Trinitate} VIII and his homilies on First John: charity turns the heart toward invisible God precisely by taking visible form in mercy and righteousness. Mysticism that ends in fascination with one's own interior state has not reached the Johannine summit.

The creature's answer is therefore both ecstatic and ordinary. To receive another person's being as gift; to will the true good rather than consume; to forgive without calling evil good; to praise without possessing; to let truth purify desire---these are finite acts, enabled and elevated by grace, in which received love returns toward unreceived Love.

The metaphysical feeling has become moral. If my existence is gift, I cannot treat another received existence as raw material for my will. If divine love creates goodness rather than feeding upon it, created love approaches its source by giving, guarding, and serving the good it did not originate.
